% Accepted replacement for the contributions paragraph:
The proposed approach combines two methodological components.
First, source identities are retained in the timing model so
that the planner associates each data-arrival contribution
with the register clock that can move it. Second, the
constructor retains multiple subtree implementations through
bounded dynamic programming (DP), allowing timing and
estimated energy to compete during construction. The resulting
circuit is evaluated through native timing analysis, matched
power simulation, and test vector leakage assessment (TVLA).
These complementary evaluations characterize the alignment
gains and their associated timing, power, and leakage-statistic
tradeoffs.

% Benchmark attribution, in Experimental Methodology:
Our evaluation circuit comprises sixteen parallel instances
of the masked PRESENT S-box implementation obtained from
the TETRIS artifact~\cite{tetris}. We retain the original
S-box RTL and assemble these instances into a 64-bit
substitution layer for physical-design evaluation.

% Add inside thebibliography, immediately after bibitem{present}:
\bibitem{tetris}
N. Sarma, T. Patidar, N. Brahamanya, and C. Karfa,
``TETRIS: Automated Design Space Exploration of
Randomness--Latency Trade-offs in Masked Hardware,''
\emph{IACR Transactions on Cryptographic Hardware and
Embedded Systems}, vol. 2026, no. 4, 2026.
Artifact: \texttt{github.com/\allowbreak nilotpolas/\allowbreak TETRIS-artifact}.
